% Co-governed and enforced under the Sovereign Integrity Protocol License (SIP License v1.1):
% https://github.com/tensorrent/ACS-Framework/blob/main/LICENSE
\documentclass[11pt]{article}
\usepackage[margin=1in]{geometry}
\usepackage{amsmath,amssymb,amsthm}
\usepackage{booktabs}
\usepackage{array}
\usepackage{enumitem}
\usepackage{xcolor}
\usepackage{hyperref}
\hypersetup{colorlinks=true,linkcolor=blue!50!black,urlcolor=blue!50!black,citecolor=blue!50!black}
\usepackage{fancyhdr}
\usepackage[protrusion=true,expansion=false]{microtype}

\pagestyle{fancy}
\fancyhf{}
\fancyhead[L]{\small Scaled Invariance of Infinity and Zero}
\fancyhead[R]{\small The Unit Is the Reference Frame}
\fancyfoot[C]{\thepage}

\theoremstyle{definition}
\newtheorem{definition}{Definition}
\theoremstyle{plain}
\newtheorem{principle}{Principle}
\newtheorem{proposition}{Proposition}
\newtheorem{corollary}{Corollary}
\newtheorem{observation}{Observation}

\newcommand{\Z}{\mathbb{Z}}
\newcommand{\R}{\mathbb{R}}
\newcommand{\Q}{\mathbb{Q}}

\title{\textbf{Scaled Invariance of Infinity and Zero:\\
The Counting Unit Is the Reference Frame}\\[4pt]
\large How ``infinitely many'' and ``none'' become two readings of one interval}
\author{B.\ Wallace \quad (TensorRent)\\
\small with computational verification by Claude (Anthropic)}
\date{\small Companion to \textit{Form/Function Relativity} (FF06g) and \textit{When a Number Lies}\ ---\ SIP License v1.1\ ---\ Seed \texttt{20260423}}

\begin{document}
\maketitle

\begin{abstract}
There are infinitely many numbers between $1$ and $2$ in the decimals, and there are none
between them in the whole numbers. Both statements are true, of the same pair $(1,2)$, and the
framework's discipline is to notice that a quantity carrying two opposite values cannot be a
property of the object it is attached to. It is a property of a \emph{pair}. We make this precise:
the count of representable points inside an interval is a reading of $(\text{interval},\,\text{unit})$,
where the unit $\delta$ is the minimum resolvable step --- a freely chosen reference frame, exactly
the role played by the surrogate ensemble in \textit{Form/Function Relativity} (FF06g) and by the
reference measure $\mathrm{d}\mu/\mathrm{d}\nu$ in the core monograph's asymmetry lemma. The labels
$\infty$ (as $\delta\to 0$) and $0$ (interior of $(1,2)$ at $\delta=1$) are the same interval read at
two units related by a scale. We isolate the invariant: the count is preserved \emph{exactly} under
the joint rescaling $(\,[a,b],\delta)\mapsto(\lambda[a,b],\lambda\delta)$ --- verified with exact
rational arithmetic over $2\times10^{5}$ random configurations, $0$ mismatches --- and it moves under
\emph{interval-only} scaling, which is why the label is relative rather than intrinsic. The only
frame-free content is the dimensionless ratio $\mathrm{width}/\delta$. The interior of $(1,2)$ carries
$0,9,99,999,9999$ points at $\delta = 1,\tfrac1{10},\tfrac1{100},\tfrac1{1000},\tfrac1{10000}$: one
interval, the whole staircase from ``empty'' to ``diverging,'' driven only by the unit. This is a
classical, machine-verified (T1) statement about resolution-relative counting. \textbf{It is not a
claim about set cardinality}: $|\R|,|\Q|,|\Z|$ remain genuinely distinct and frame-independent, and
Cantor is nowhere touched. The observation is the reframing; the scope boundary is the point.
\end{abstract}

\section{The observation}\label{sec:obs}

Ask how many numbers lie between $1$ and $2$. In the decimals the answer is ``infinitely many'';
in the whole numbers the answer is ``none.'' The two answers are not in tension --- they are answers
to two different questions wearing the same words. What changed between them is not the interval
$[1,2]$, which is fixed, but the \emph{unit}: the smallest step you are permitted to resolve. The
program's first move, made already in \textit{When a Number Lies} and \textit{Form/Function
Relativity}, is to refuse to attach a swinging label to a fixed object. If the count of interior
points takes the value $\infty$ under one convention and $0$ under another, then the count is not a
property of $[1,2]$. It is a property of the \emph{pair} $\big([1,2],\,\delta\big)$, and $\delta$ ---
the resolution unit --- is a chosen reference frame.

\section{The unit as reference frame}\label{sec:frame}

\begin{definition}[Resolution unit and its count]\label{def:count}
Fix a unit $\delta>0$: the minimum resolvable step, so that the only representable points are the
lattice $\delta\Z=\{m\delta : m\in\Z\}$. For an interval $[a,b]$ define the closed count
\[
  N_\delta[a,b]\;=\;\#\big(\delta\Z\cap[a,b]\big)\;=\;\big\lfloor b/\delta\big\rfloor-\big\lceil a/\delta\big\rceil+1
  \quad(\text{clamped at }0),
\]
and the interior count $N_\delta^{\circ}(a,b)=\#\big(\delta\Z\cap(a,b)\big)$. The ``decimals'' limit is
$\delta\to 0$; the ``whole numbers'' frame is $\delta=1$.
\end{definition}

The count is now explicitly a two-argument object $N_\delta[a,b]$, and the content of this note is
that the second argument is load-bearing and that its dependence has an exact symmetry.

\begin{principle}[Scaled invariance of the counting label]\label{prin:scaled}
The count of representable points in an interval is a coordinate on the pair
$(\text{interval}\times\text{unit})$, not a property of the interval. It is \emph{invariant} under the
joint rescaling $(\,[a,b],\delta)\mapsto(\lambda[a,b],\lambda\delta)$ for every $\lambda>0$, and it
\emph{varies} under scaling either factor alone. The labels ``$\infty$'' and ``$0$'' for the interior
of a fixed interval are therefore relative to the unit, and are related to one another by a change of
scale --- infinity and zero are the two ends of one staircase in $\delta$.
\end{principle}

\begin{proposition}[Exact joint-scaling invariance]\label{prop:invariance}
For all $a\le b$, $\delta>0$, $\lambda>0$,
\[
  N_\delta[a,b]\;=\;N_{\lambda\delta}\big[\lambda a,\lambda b\big].
\]
\end{proposition}
\begin{proof}
The map $x\mapsto\lambda x$ is a bijection $\delta\Z\to(\lambda\delta)\Z$ carrying $[a,b]$ onto
$[\lambda a,\lambda b]$: $x\in\delta\Z\cap[a,b]\iff \lambda x\in(\lambda\delta)\Z\cap[\lambda a,\lambda b]$.
A bijection preserves cardinality. Equivalently, both counts equal
$\lfloor b/\delta\rfloor-\lceil a/\delta\rceil+1$ because $(\lambda b)/(\lambda\delta)=b/\delta$ and
$(\lambda a)/(\lambda\delta)=a/\delta$. \qed
\end{proof}

\begin{corollary}[The invariant is dimensionless]\label{cor:dimless}
$N_\delta[a,b]$ depends only on the ratios $a/\delta$ and $b/\delta$. The raw width and the raw unit
are gauge; only $\mathrm{width}/\delta$ is physical. In particular the interior count between two
consecutive units is $0$ for \emph{every} $\delta$ --- ``adjacent'' means ``one unit apart,'' a
statement about $\delta$, never about the two endpoints.
\end{corollary}

Proposition~\ref{prop:invariance} is the precise form of ``infinity and zero are only a relative
scaled invariance.'' Scale the interval and the unit together and nothing moves; the count is a rigid
invariant. Scale one without the other and the label slides continuously from $0$ to $\infty$. There
is no absolute count of ``the numbers between $1$ and $2$''; there is only the count at a stated
resolution, and a symmetry that says which reparametrizations leave it fixed.

\section{Measurement}\label{sec:meas}

The counting facts are exact integers; we verify them with exact rational arithmetic (Python
\texttt{fractions}, no floating error), seed \texttt{20260423}. The driver is
\texttt{code/notes\_verification/test\_scaled\_invariance.py}.

\paragraph{A. Unit-relativity of the label.} The interior of the single fixed interval $(1,2)$, read
at shrinking units:

\begin{center}
\renewcommand{\arraystretch}{1.3}
\begin{tabular}{@{}lccccc@{}}
\toprule
unit $\delta$ & $1$ & $\tfrac{1}{10}$ & $\tfrac{1}{100}$ & $\tfrac{1}{1000}$ & $\tfrac{1}{10000}$\\
\midrule
$N_\delta^{\circ}(1,2)$ & $0$ & $9$ & $99$ & $999$ & $9999$\\
label & empty & \multicolumn{4}{c}{$\longrightarrow$ diverging $(\to\infty)$}\\
\bottomrule
\end{tabular}
\end{center}

\noindent One interval; the entire span from ``nothing between $1$ and $2$'' to ``infinitely many
between $1$ and $2$'' is produced by the unit alone. Principle~\ref{prin:scaled} made numerical.

\paragraph{B. Joint-scaling invariance (the symmetry).} Over $2\times10^{5}$ random configurations
$([a,b],\delta,\lambda)$ with $a,b,\delta,\lambda$ rational, comparing $N_\delta[a,b]$ against
$N_{\lambda\delta}[\lambda a,\lambda b]$: \textbf{$0$ mismatches} (largest count encountered $45{,}601$).
Proposition~\ref{prop:invariance} holds exactly, as a bijection must.

\paragraph{C. Interval-only scaling breaks it (why it is ``relative'').} Holding $\delta=1$ and
stretching the interval:
\[
  N_1^{\circ}(1,2)=0,\quad N_1^{\circ}(10,20)=9,\quad N_1^{\circ}(100,200)=99,\quad N_1^{\circ}(1000,2000)=999.
\]
The $\times 10$ image of $[1,2]$ gains $9$ interior points \emph{because the unit did not scale with
it}. The adjacency of $1$ and $2$ is the statement $b-a=\delta$, a fact about the frame.

\paragraph{D. The invariant is dimensionless.} Over $5\times10^{4}$ random ratio-preserving
reparametrizations, $N_\delta[a,b]=N_{\lambda\delta}[\lambda a,\lambda b]$ in \textbf{$50{,}000/50{,}000$}
cases, confirming Corollary~\ref{cor:dimless}: the count is a function of $a/\delta,b/\delta$ alone.

All four are machine-verified (T1); the script exits nonzero on any mismatch.

\section{Why this is the same object as the framework's other relativities}\label{sec:same}

This is not a new mechanism; it is the framework's standing one, applied to counting.

\textbf{The unit is the reference measure.} \textit{Form/Function Relativity} (FF06g) showed the
form/function verdict to be a coordinate on (witness $\times$ frame), with the frame a chosen surrogate
ensemble $\nu$ --- the reference measure $\mathrm{d}\mu/\mathrm{d}\nu$ of the core monograph's
BCH--transfer-entropy lemma. Here the ``frame'' is the resolution unit $\delta$ and the ``label'' is
the cardinality reading $\{0,\dots,\infty\}$ rather than $\{$form, function$\}$. In both cases a scalar
that looks intrinsic is exhibited as a two-argument object, and the second argument --- the frame --- is
free. Choosing it differently relabels; that choice is the perspective.

\textbf{The scalar was a shadow.} \textit{When a Number Lies} isolates the pattern that a single scalar
is often a lossy projection of an underlying relation, and that the relation, not the scalar, is the
faithful object. ``The number of points between $1$ and $2$'' is such a shadow: the faithful object is
the relation $\mathrm{width}/\delta$ between the interval and the unit, and only that dimensionless
ratio survives a change of frame (Corollary~\ref{cor:dimless}). ``$\infty$'' and ``$0$'' are what the
shadow reads when the unit is sent to $0$ or set equal to the width.

\textbf{Form/function is frame-relative; so is cardinality-under-resolution.} The parent program's
thesis is that form and function are generated by the asymmetry of a pair and carried by neither member
alone. Principle~\ref{prin:scaled} is the counting instance: ``how many'' is generated by the asymmetry
between an interval and a unit, and is carried by neither the interval nor the unit alone.

\section{What this establishes, and where it stops}\label{sec:scope}

\textbf{Establishes.} For a fixed interval, the count of representable points is frame-relative in the
exact sense of Definition~\ref{def:count}: it ranges over $0,9,99,999,9999,\dots$ on the single interval
$(1,2)$ as the unit shrinks (measurement A), it is preserved \emph{exactly} under joint rescaling of
interval and unit (Proposition~\ref{prop:invariance}, $0/2\!\times\!10^5$ mismatches), it changes under
interval-only scaling (measurement C), and its only frame-free content is the dimensionless ratio
$\mathrm{width}/\delta$ (Corollary~\ref{cor:dimless}, $50000/50000$). ``Infinity'' and ``zero'' are the
$\delta\to0$ and $\delta=\mathrm{width}$ ends of one staircase, related by a scale.

\textbf{Does not establish --- the load-bearing boundary.} This is a statement about the count of
points \emph{representable at a resolution}, i.e.\ the lattice $\delta\Z$. It is \textbf{not} a statement
about set cardinality. The claim ``there are infinitely many reals between $1$ and $2$'' is true and
frame-\emph{independent} as a statement about $\R$: $|\,[1,2]\cap\R\,|=\mathfrak{c}$ and
$|\,[1,2]\cap\Q\,|=\aleph_0$ regardless of any unit, because those sets contain their points with no
resolution imposed. Cantor's theorem, the distinctness of $\aleph_0$ and $\mathfrak c$, and the density
of $\Q$ in $\R$ are untouched and uncontested here. The relativity is precisely the difference between
\emph{``how many points exist in the set''} (cardinality; frame-free; the reals are genuinely
uncountable) and \emph{``how many points are resolvable under a minimum step $\delta$''} (a count;
frame-relative; what this note is about). Conflating the two would be the overclaim, and the boundary
between them \emph{is} the result: the whole-number constraint does not make the reals disappear --- it
selects a coarser frame in which the interior of $(1,2)$ is empty, and Proposition~\ref{prop:invariance}
says exactly how that frame relates to the finer ones.

\textbf{Two honest edges.} (i) $N_\delta^{\circ}(1,2)=k-1$ at $\delta=1/k$ grows without bound but is
finite at every $\delta>0$; ``$\infty$'' is the limit $\delta\to0$, not a value at any admissible unit,
and the limit is exactly where the resolution frame is abandoned and the set-cardinality question (a
different question) takes over. (ii) The invariance is a discrete lattice fact, not a measure-theoretic
one; the continuous analogue --- Lebesgue length $b-a$ scaling by $\lambda$ under $x\mapsto\lambda x$ ---
is the $\delta\to0$ shadow of Proposition~\ref{prop:invariance} with the joint rescaling of the unit
made explicit rather than hidden. Both edges are recorded so the staircase is not over-read.

\vspace{1em}
\noindent\rule{\linewidth}{0.4pt}\\
\small This companion accompanies \textit{Form/Function Relativity} (FF06g) and \textit{When a Number
Lies}. All counts are exact-integer and reproduced standalone via
\texttt{code/notes\_verification/test\_scaled\_invariance.py} at seed \texttt{20260423}. SIP License v1.1.

\end{document}
